% ============================================================
% レポートの書き方（手引書の前付け，注意事項の次に入れる）
% 参考にした資料:
%   基礎ゼミ実験配布テキスト.pdf（電気系基礎実験A・B 実験の手引き）
%   基礎ゼミ第2回資料.pdf（基礎ゼミナール 実験ガイダンス：実験報告書の作成方法）
%   レポートの書き方.pdf（他大学の実験上の諸注意・報告書の書き方）
% 図・表・式のラベルは report- で始める
% ============================================================

\introtitle{レポートの書き方}

レポートは，実験の方法・過程・結果・考察が読む人に正しく伝わり，正当に評価されるように，十分な時間をかけて丁寧に作成してください。ここでは，レポートを書くときの基本的な決まりをまとめます。

% ------------------------------------------------------------
\section{全般的な注意}

\begin{itemize}
  \item レポートは授業時間外に作成します。図や表は，実験が終わったら記憶が新しいうちに速やかに作成してください。
  \item レポートは\wave{他の人に読まれるものであるという意識をもって，丁寧に作成してください。}
  \item 考察は必ず自分自身で考えて書いてください。レポートは自分の著作物です。他人のレポートや文章をそのまま写すことは認められません。
  \item 生成系AIのみを用いて作成したレポートを提出してはいけません。\wave{生成系AIによる出力をそのまま提出したことが判明した場合は，単位を認定しない，または認定を取り消すことがあります。}
  \item レポートの文章は「〜である」「〜だ」の調子（である調）で書きます。手引書は「〜します」「〜してください」と書いていますが，レポートとしては不適切です。
  \item 用紙はA4のレポート用紙とします。報告書には表紙を付けてください（表紙の書き方は\pageref{page:report-cover}ページの表紙例を参照）。表紙以外の各ページにはページ番号を入れます。
\end{itemize}

% ------------------------------------------------------------
\section{レポートの構成}

レポートは，表紙と本文からなります。表紙には，実験題目，実験日，報告書の提出日，班，報告者の学籍番号と氏名，共同実験者の学籍番号と氏名を記入します（\pageref{page:report-cover}ページの表紙例を参照）。再提出するときは，報告書再提出日も記入します。

本文は，\cref{tab:report-structure}のような章に分けて，順序よく書きます。必要に応じて節（2.1節，2.1.1節など）に分けると，内容がわかりやすくなります。図番号・表番号は通し番号にし，各図・表にはその説明（キャプション）を入れます。

\begin{center}
\begin{minipage}{\linewidth}% 見出しと表を分けない
  \centering
  \captionof{table}{レポート本文の構成}
  \label{tab:report-structure}
  \renewcommand{\arraystretch}{1.3}
  \begin{tabular}{@{}l>{\raggedright\arraybackslash}p{0.68\linewidth}@{}}
    \toprule
    章 & 書く内容 \\
    \midrule
    1.~目的     & 実験で何を明らかにするのか，何を確かめるのかを簡潔に書く。 \\
    2.~原理     & 実験の基礎となる知識や理論を要領よくまとめる。必要な式や回路図も示す。 \\
    3.~使用器具 & 実験に使用した器具を，後から特定できるように書く（品名，メーカー名，型番，定格・性能など）。 \\
    4.~実験方法 & 実験の手順，実験回路，設定した値（パラメータ）などを具体的に書く。 \\
    5.~実験結果 & 測定値やそこから計算した値を，図や表にまとめて示し，説明する。 \\
    6.~考察     & 実験結果から何がいえるのか，なぜそのような結果になったのかを論じる。 \\
    7.~参考文献 & 考察などで参考にした文献を，本文中の番号と対応させて示す。 \\
    \bottomrule
  \end{tabular}
\end{minipage}
\end{center}

\begin{hosoku}[使用器具の書き方]
使用器具は，後でレポートを読んだ人が同じ器具を特定できるように書きます。例えば直流電源であれば，
\begin{center}
  直流定電圧電源（メーカー：○○，型番：○○，定格：\SI{35}{V}，\SI{0.5}{A}）
\end{center}
のように書きます。定格（出力できる最大の電圧・電流や，測定できる最大の電圧・電流）や性能（測定精度など）を記録しておくと，結果がおかしいときに原因を調べる手がかりになります。
\end{hosoku}

\begin{hosoku}[実験結果の書き方]
実験結果は，実験ごとに「どのような実験の結果か（実験の条件）」「どの図・表に示したか」「図・表から何が読み取れるか」をはっきり書きます。測定値をただ並べるのではなく，図や表に整理して示します。結果のグラフや表は，レポートの最後にまとめるのではなく，実験結果の章の中に入れます。
\end{hosoku}

% ------------------------------------------------------------
\section{図の書き方}

図は，その図だけを見ても内容がわかるように（図だけで独り歩きできるように）書きます。グラフも「図」として扱い，「グラフ1」のような別の番号の付け方をしてはいけません。グラフの例を\cref{fig:report-graph}に示します。

\begin{figure}[htbp]
  \centering
  \begin{tikzpicture}
    \begin{axis}[
        width=0.62\linewidth, height=0.42\linewidth,
        xmin=0, xmax=10, ymin=0, ymax=10,
        xtick={0,2,...,10}, ytick={0,2,...,10},
        xlabel={電圧 $V$ (V)}, ylabel={電流 $I$ (mA)},
        tick align=inside, axis on top,
        legend pos=north west, legend cell align=left,
        legend style={font=\footnotesize, draw=none, fill=none},
      ]
      \addplot[only marks, mark=o, mark size=2.2pt, thick]
        coordinates {(1,0.98) (2,2.03) (3,2.96) (4,4.02) (5,5.01) (6,5.97) (7,7.04) (8,7.95) (9,9.02)};
      \addlegendentry{測定値}
      \addplot[domain=0:9.8, samples=2, black]{x};
      \addlegendentry{近似直線（最小二乗法）}
    \end{axis}
  \end{tikzpicture}
  \caption{抵抗器の電圧と電流の関係（作成例）}
  \label{fig:report-graph}
  {\footnotesize 供試物：金属皮膜抵抗器（\SI{1}{k\ohm}）}
\end{figure}

\begin{itemize}
  \item 図番号（通し番号）と図の説明（キャプション）は，\important{図の下}に書きます。「○○の××依存性」「○○−××特性」と書くときは，○○が縦軸，××が横軸の物理量です。
  \item 縦軸・横軸には，物理量の名前と単位を，軸の\important{中央}に「電圧 $V$ (V)」のように書きます。軸の先に矢印は付けません。
  \item 目盛は等間隔に入れ，軸の両端に目盛の数字がくるようにします。目盛は枠の内側に入れます。軸の範囲はデータのある範囲に合わせます。測定値の数字を軸に書き込んではいけません。
  \item 測定値は，○，●，△，×などの印で，適度な大きさではっきりとプロットします。
  \item 測定値を1点ずつ線でつなぐのではなく，最も確からしい1本のなめらかな線（理論的に直線になる場合は最小二乗法で求めた直線）を引きます。
  \item 1つの図に複数の結果を描くときは，印や線の種類を変え，それぞれが何を表すか（凡例）を書き入れます。
  \item 実験条件や供試物は，図の中か，図の説明の下に書きます。
  \item 図はその図について説明した文の後に置き，本文中で必ず「図○に示すように…」と説明します。
\end{itemize}

% ------------------------------------------------------------
\section{表の書き方}

表の例を\cref{tab:report-example}に示します。

\begin{table}[htbp]
  \centering
  \caption{抵抗器の電圧と電流の関係（作成例）}
  \label{tab:report-example}
  \begin{tabular}{cccc}
    \toprule
    設定電圧 $V_\mathrm{S}$ (V) & 測定電圧 $V$ (V) & 電流 $I$ (mA) & 抵抗 $R$ (k$\Omega$) \\
    \midrule
    1.0 & 0.99 & 0.98 & 1.01 \\
    2.0 & 2.01 & 2.03 & 0.990 \\
    3.0 & 2.98 & 2.96 & 1.01 \\
    \bottomrule
  \end{tabular}
\end{table}

\begin{itemize}
  \item 表番号（通し番号）と表の説明は，\important{表の上}に書きます。
  \item 各列の見出しには，物理量の名前と単位を「電流 $I$ (mA)」のように書きます。
  \item 数値は有効数字を考えて書きます（\cref{sec:report-figures}を参照）。\SI{0.001}{A}のような値は，接頭語を使って\SI{1}{mA}のように書くと読みやすくなります。
  \item 表はページをまたいで書かないようにします。
  \item 表も，本文中で必ず説明します。
\end{itemize}

% ------------------------------------------------------------
\section{数値と単位の書き方}
\label{sec:report-figures}

\begin{hosoku}[有効数字]
測定値には必ず誤差があり，意味のある桁数（\important{有効数字}）は限られています。例えば，1.60は有効数字3桁（$1.60\pm0.005$），1.6は有効数字2桁（$1.6\pm0.05$）を表します。1600のように0で終わる数は，0が単なる位取りなのか意味のある数字なのかわからないので，$1.6\times10^{3}$（2桁），$1.60\times10^{3}$（3桁）のように書きます。電卓で計算した値をすべての桁まで書くのは意味がありません。
\end{hosoku}

\begin{hosoku}[計算と有効数字]
\begin{itemize}
  \item \important{かけ算・わり算}では，結果の有効数字の桁数を，もとの数のうち桁数の少ない方に合わせます。\\
    例：$1.3\times2.12=2.756$ → $2.8$（有効数字2桁）
  \item \important{足し算・引き算}では，もとの数のうち最も上の位までしか意味のない数に合わせて，結果の位をそろえます。\\
    例：$23+1.5=24.5$ → $25$（一の位まで）
  \item 近い値どうしの引き算では，有効数字の桁数が大きく減ります（桁落ち）。例：$113.9-110.7=3.2$（2桁）
\end{itemize}
\end{hosoku}

\begin{hosoku}[単位]
数値には必ず単位を付けます（\SI{500}{V}，\SI{213}{mA}）。グラフの軸や表の見出しのように，物理量の記号に単位を添えるときは「電圧 $V$ (V)」のように単位を括弧に入れます。
\end{hosoku}

% ------------------------------------------------------------
\section{考察の書き方}

考察では，実験結果が何を示しているのか，結果を総合して何がいえるのかを論じます。
\begin{itemize}
  \item 感想ではなく，実験から得られた\important{結論}を書きます。
  \item 実験結果の説明，結果からわかること，なぜそのような結果になったのかの理由を，原理（理論）と比べながら明確に書きます。
  \item 理論と合わない結果が得られたときは，その原因を具体的に考えます。合わない結果をすべて「測定誤差のため」で片付けるのは考察になりません。
  \item 実験の問題点や改善点についても考えます。
  \item 考察はいくつかの小項目に分けて書くと読みやすくなります。各テーマの「考察と課題」は，考察を書くための方向性を示したものです。課題に答えるだけでなく，自分で結果を検討してください。
\end{itemize}

\begin{hosoku}[考察の例]
\textbf{よくない例：}抵抗値は理論値とだいたい一致した。うまく測定できてよかった。

\textbf{よい例：}\cref{fig:report-graph}より，電流 $I$ は電圧 $V$ にほぼ比例しており，最小二乗法で求めた直線の傾きから抵抗値は\SI{1.00}{k\ohm}と求められた。これは抵抗器の表示値\SI{1}{k\ohm}（許容差\SI{\pm1}{\percent}）の範囲内であり，オームの法則が成り立つことが確かめられた。
\end{hosoku}

% ------------------------------------------------------------
\section{参考文献の書き方}

考察などで参考にした文献は，本文中の参考にした箇所に通し番号を付け，レポートの最後にまとめて示します。
\begin{center}
\begin{minipage}{0.9\linewidth}
  本文中：…銅の室温における抵抗率は $1.68\times10^{-8}\ \Omega\,\mathrm{m}$ であることが知られている\,[1]。…\\[0.5\baselineskip]
  参考文献：\\
  {[1]} 著者名：書名，ページ，出版社（発行年）\\
  {[2]} 著者名：論文名，雑誌名，巻，ページ（発行年）\\
  {[3]} 著者名（運営者名）：ページの題名，URL（閲覧日）
\end{minipage}
\end{center}
Webページを参考文献とする場合，個人のブログなど信頼性が確認できない情報はなるべく避け，公的機関やメーカー，学会などの信頼できる情報源を選んでください。

% ------------------------------------------------------------
% 表紙例（1ページ）。基礎ゼミ実験配布テキスト.pdf の p.8「実験報告書の表紙例」をもとに作成
% 枠は A4 と同じ縦横比（120 mm × 169.7 mm = 1 : 1.414）
\clearpage
\phantomsection
\label{page:report-cover}
\noindent{\sffamily\gtfamily\bfseries\large 実験報告書の表紙例}

\vspace{1.0\baselineskip}
\begin{center}
\begin{tikzpicture}[
    lb/.style={anchor=west, font=\normalsize},
    ln/.style={line width=0.5pt},
  ]
  \draw[line width=0.6pt] (0,0) rectangle (120mm,169.7mm);
  \node[font=\sffamily\gtfamily\LARGE] at (60mm,158mm) {実験報告書};
  % 実験題目
  \node[lb] at (12mm,139mm) {実験題目};
  \draw[ln] (20mm,128mm) -- (108mm,128mm);
  % 日付
  \foreach \y/\t in {114/実験開始日, 105/実験終了日, 96/報告書提出日, 87/{\llap{*}報告書再提出日}} {
    \node[lb] at (32mm,\y mm) {\t};
    \node[lb] at (66mm,\y mm) {令和\hspace{2.2em}年\hspace{2.2em}月\hspace{2.2em}日};
  }
  \node[lb, font=\footnotesize] at (32mm,80mm) {（*必要に応じて記入する）};
  % 班
  \draw[ln] (14mm,66mm) -- (32mm,66mm);
  \node[lb] at (32.5mm,67.2mm) {班};
  % 報告者
  \node[lb] at (12mm,54mm) {報告者};
  \node[lb] at (16mm,47mm) {学籍番号};
  \node[lb] at (60mm,47mm) {氏名};
  \draw[ln] (14mm,39mm) -- (48mm,39mm);
  \draw[ln] (60mm,39mm) -- (108mm,39mm);
  % 共同実験者
  \node[lb] at (12mm,30mm) {共同実験者};
  \node[lb] at (16mm,23mm) {学籍番号};
  \node[lb] at (60mm,23mm) {氏名};
  \foreach \y in {16,10} {
    \draw[ln] (14mm,\y mm) -- (48mm,\y mm);
    \draw[ln] (60mm,\y mm) -- (108mm,\y mm);
  }
\end{tikzpicture}
\end{center}
